\begin{theorem}[最小非平凡模数已经越过平方根门槛]\label{thm:xi-real-input-40-collision-minimal-modulus-two-witness}
沿用定理 \ref{thm:xi-real-input-40-collision-prime-twist-p2-threshold} 与
\ref{thm:xi-real-input-40-collision-mod2-trace-three-spectrum} 的记号。对每个 \(m\ge 2\)，定义
\[
\omega_m:=e^{2\pi i/m},
\qquad
\rho_{m,j}:=\rho\!\bigl(W(\omega_m^j)\bigr),
\qquad
\theta_{m,j}:=\frac{\log \rho_{m,j}}{\log\lambda},
\]
\[
\theta_m:=\max_{1\le j\le m-1}\theta_{m,j},
\qquad
\lambda=\varphi^2.
\]
则
\[
\theta_2=\theta_{2,1}>\frac12.
\]
特别地，\(m=2\) 已经是最小非平凡模数上的平方根门槛失效证人。
\end{theorem}
\begin{proof}
当 \(m=2\) 时，唯一非平凡单位根为
\[
\omega_2=-1,
\]
故
\[
\theta_2=\theta_{2,1}=\frac{\log\rho_-}{\log\lambda},
\]
其中 \(\rho_-\) 是定理 \ref{thm:xi-real-input-40-collision-mod2-trace-three-spectrum} 中三次方程
\[
f(x):=x^3-2x^2-1=0
\]
的最大实根。

现直接在 \(x=\varphi\) 处检验。由
\[
\varphi^2=\varphi+1,
\qquad
\varphi^3=2\varphi+1,
\]
得到
\[
f(\varphi)
=
\varphi^3-2\varphi^2-1
=
(2\varphi+1)-2(\varphi+1)-1
=
-2<0.
\]
另一方面，
\[
f'(x)=3x^2-4x=x(3x-4)>0
\qquad (x\ge \varphi),
\]
因为 \(\varphi>4/3\)。故 \(f\) 在区间 \([\varphi,\infty)\) 上严格递增。既然
\[
f(\varphi)<0,
\qquad
f(x)\to +\infty\quad(x\to+\infty),
\]
其最大实根必满足
\[
\rho_->\varphi=\sqrt{\lambda}.
\]
取对数并除以 \(\log\lambda>0\)，便得到
\[
\theta_2=\frac{\log\rho_-}{\log\lambda}
>
\frac{\log\sqrt{\lambda}}{\log\lambda}
=
\frac12.
\]
最后，由于 \(m=2\) 是存在非平凡单位根扭曲的最小模数，结论成立。证毕。
\end{proof}

\begin{theorem}[最坏 Dirichlet 扭曲指数逼近 \texorpdfstring{$1$}{1} 并排斥任何统一次指数节约]\label{thm:xi-real-input-40-collision-worst-twist-exponent-tends-to-one}
沿用定理 \ref{thm:xi-real-input-40-collision-minimal-modulus-two-witness} 的记号，并定义
\[
\rho_m:=\max_{1\le j\le m-1}\rho_{m,j}.
\]
则当 \(m\to\infty\) 时有显式下界
\[
\boxed{\;
\theta_m
\ge
1-
\frac{(6\sqrt5-5)(2\pi)^2}{250\log(\varphi^2)}\,\frac1{m^2}
+
\frac{(7+24\sqrt5)(2\pi)^4}{75000\log(\varphi^2)}\,\frac1{m^4}
+
O(m^{-6}).\;}
\]
特别地，
\[
\theta_m\longrightarrow 1
\qquad (m\to\infty).
\]
因此不存在任何统一常数 \(\sigma<1\) 使得对全部 \(m\ge 2\) 都有
\[
\rho_m\le \lambda^\sigma.
\]
等价地，对任意 \(\sigma<1\)，都有充分大的全部模数 \(m\) 满足
\[
\theta_m>\sigma.
\]
\end{theorem}
\begin{proof}
取最小非平凡角
\[
t_m:=\frac{2\pi}{m}.
\]
则
\[
\omega_m=e^{it_m}
\]
是 \(m\)-次单位根，因此
\[
\rho_m\ge \rho_{m,1}=\rho\!\bigl(W(\omega_m)\bigr).
\]
由推论 \ref{cor:real-input-40-collision-twist-fourth}，
\[
\log\frac{\rho(W(e^{it}))}{\lambda}
=
-\frac{6\sqrt5-5}{250}\,t^2
+\frac{7+24\sqrt5}{75000}\,t^4
+O(t^6)
\qquad (t\to 0).
\]
代入 \(t=t_m=2\pi/m\)，得到
\[
\log\frac{\rho_{m,1}}{\lambda}
=
-\frac{(6\sqrt5-5)(2\pi)^2}{250}\,\frac1{m^2}
+\frac{(7+24\sqrt5)(2\pi)^4}{75000}\,\frac1{m^4}
+O(m^{-6}).
\]
两边除以 \(\log\lambda\)，并利用 \(\theta_m\ge \theta_{m,1}\)，即得所示下界。

另一方面，对任意单位根 \(\omega_m^j\)，都有
\[
|W(\omega_m^j)|=W(1)
\]
逐项按模取值成立。故由谱半径比较，
\[
\rho_{m,j}\le \rho(W(1))=\lambda,
\]
从而
\[
\theta_m\le 1.
\]
于是由下界与上界夹逼可知
\[
\theta_m\to 1.
\]

若存在某个 \(\sigma<1\) 使得 \(\rho_m\le \lambda^\sigma\) 对所有 \(m\ge 2\) 都成立，则必有
\[
\theta_m=\frac{\log\rho_m}{\log\lambda}\le \sigma<1
\qquad(\forall m\ge 2),
\]
这与 \(\theta_m\to 1\) 矛盾。故不存在这样的统一 \(\sigma\)。证毕。
\end{proof}

